\section{Conclusion}
\label{sec:conclusion}

We have shown that the $\Sigma_1$ theory of non-structural subtyping is
decidable for one type constructor of any arity and variance with $\Bot$ and
$\Top$, over arbitrary, regular and finite trees, and for a covariant
constructor with $\Bot$ only. Equivalently, it is decidable whether a
constraint entails a disjunction of inequalities. For a binary covariant
constructor, where the $\exists\forall\exists$ fragment is undecidable, this
locates the boundary of decidability in the quantifier hierarchy between the
first and the third level.

The proof reduces a countermodel with $k$ failed inequalities to a tuple of
$k$ words that avoids forbidden pairs of states wherever the rest of one
word is a prefix of the rest of another, and bounds the length of a shortest
such tuple. The two halves of the bound treat two different phenomena. Words
that can be continued forever must be prevented from running into each
other, which leads to periodic continuations, difference constraints between
their lengths, and the compression of their periodic stretches. After the tails are
bounded, the cover turns the hole into a fresh tuple, which is shortened by
deleting the same repeated blocks from all its words along a hierarchy of
markers. The resulting bound is elementary, and
the whole argument, from the semantics of trees to the closed form of the
bound, is checked in Lean.
